% Options for packages loaded elsewhere
\PassOptionsToPackage{unicode}{hyperref}
\PassOptionsToPackage{hyphens}{url}
\documentclass[
]{article}
\usepackage{xcolor}
\usepackage{amsmath,amssymb}
\setcounter{secnumdepth}{-\maxdimen} % remove section numbering
\usepackage{iftex}
\ifPDFTeX
  \usepackage[T1]{fontenc}
  \usepackage[utf8]{inputenc}
  \usepackage{textcomp} % provide euro and other symbols
\else % if luatex or xetex
  \usepackage{unicode-math} % this also loads fontspec
  \defaultfontfeatures{Scale=MatchLowercase}
  \defaultfontfeatures[\rmfamily]{Ligatures=TeX,Scale=1}
\fi
\usepackage{lmodern}
\ifPDFTeX\else
  % xetex/luatex font selection
\fi
% Use upquote if available, for straight quotes in verbatim environments
\IfFileExists{upquote.sty}{\usepackage{upquote}}{}
\IfFileExists{microtype.sty}{% use microtype if available
  \usepackage[]{microtype}
  \UseMicrotypeSet[protrusion]{basicmath} % disable protrusion for tt fonts
}{}
\makeatletter
\@ifundefined{KOMAClassName}{% if non-KOMA class
  \IfFileExists{parskip.sty}{%
    \usepackage{parskip}
  }{% else
    \setlength{\parindent}{0pt}
    \setlength{\parskip}{6pt plus 2pt minus 1pt}}
}{% if KOMA class
  \KOMAoptions{parskip=half}}
\makeatother
\usepackage{color}
\usepackage{fancyvrb}
\newcommand{\VerbBar}{|}
\newcommand{\VERB}{\Verb[commandchars=\\\{\}]}
\DefineVerbatimEnvironment{Highlighting}{Verbatim}{commandchars=\\\{\}}
% Add ',fontsize=\small' for more characters per line
\newenvironment{Shaded}{}{}
\newcommand{\AlertTok}[1]{\textcolor[rgb]{1.00,0.00,0.00}{\textbf{#1}}}
\newcommand{\AnnotationTok}[1]{\textcolor[rgb]{0.38,0.63,0.69}{\textbf{\textit{#1}}}}
\newcommand{\AttributeTok}[1]{\textcolor[rgb]{0.49,0.56,0.16}{#1}}
\newcommand{\BaseNTok}[1]{\textcolor[rgb]{0.25,0.63,0.44}{#1}}
\newcommand{\BuiltInTok}[1]{\textcolor[rgb]{0.00,0.50,0.00}{#1}}
\newcommand{\CharTok}[1]{\textcolor[rgb]{0.25,0.44,0.63}{#1}}
\newcommand{\CommentTok}[1]{\textcolor[rgb]{0.38,0.63,0.69}{\textit{#1}}}
\newcommand{\CommentVarTok}[1]{\textcolor[rgb]{0.38,0.63,0.69}{\textbf{\textit{#1}}}}
\newcommand{\ConstantTok}[1]{\textcolor[rgb]{0.53,0.00,0.00}{#1}}
\newcommand{\ControlFlowTok}[1]{\textcolor[rgb]{0.00,0.44,0.13}{\textbf{#1}}}
\newcommand{\DataTypeTok}[1]{\textcolor[rgb]{0.56,0.13,0.00}{#1}}
\newcommand{\DecValTok}[1]{\textcolor[rgb]{0.25,0.63,0.44}{#1}}
\newcommand{\DocumentationTok}[1]{\textcolor[rgb]{0.73,0.13,0.13}{\textit{#1}}}
\newcommand{\ErrorTok}[1]{\textcolor[rgb]{1.00,0.00,0.00}{\textbf{#1}}}
\newcommand{\ExtensionTok}[1]{#1}
\newcommand{\FloatTok}[1]{\textcolor[rgb]{0.25,0.63,0.44}{#1}}
\newcommand{\FunctionTok}[1]{\textcolor[rgb]{0.02,0.16,0.49}{#1}}
\newcommand{\ImportTok}[1]{\textcolor[rgb]{0.00,0.50,0.00}{\textbf{#1}}}
\newcommand{\InformationTok}[1]{\textcolor[rgb]{0.38,0.63,0.69}{\textbf{\textit{#1}}}}
\newcommand{\KeywordTok}[1]{\textcolor[rgb]{0.00,0.44,0.13}{\textbf{#1}}}
\newcommand{\NormalTok}[1]{#1}
\newcommand{\OperatorTok}[1]{\textcolor[rgb]{0.40,0.40,0.40}{#1}}
\newcommand{\OtherTok}[1]{\textcolor[rgb]{0.00,0.44,0.13}{#1}}
\newcommand{\PreprocessorTok}[1]{\textcolor[rgb]{0.74,0.48,0.00}{#1}}
\newcommand{\RegionMarkerTok}[1]{#1}
\newcommand{\SpecialCharTok}[1]{\textcolor[rgb]{0.25,0.44,0.63}{#1}}
\newcommand{\SpecialStringTok}[1]{\textcolor[rgb]{0.73,0.40,0.53}{#1}}
\newcommand{\StringTok}[1]{\textcolor[rgb]{0.25,0.44,0.63}{#1}}
\newcommand{\VariableTok}[1]{\textcolor[rgb]{0.10,0.09,0.49}{#1}}
\newcommand{\VerbatimStringTok}[1]{\textcolor[rgb]{0.25,0.44,0.63}{#1}}
\newcommand{\WarningTok}[1]{\textcolor[rgb]{0.38,0.63,0.69}{\textbf{\textit{#1}}}}
\setlength{\emergencystretch}{3em} % prevent overfull lines
\providecommand{\tightlist}{%
  \setlength{\itemsep}{0pt}\setlength{\parskip}{0pt}}
\usepackage{bookmark}
\IfFileExists{xurl.sty}{\usepackage{xurl}}{} % add URL line breaks if available
\urlstyle{same}
\hypersetup{
  hidelinks,
  pdfcreator={LaTeX via pandoc}}

\author{}
\date{}

\begin{document}

tutti gli appunti prima li ho perduti

erano la funzione map()

e... ???

Introduciamo il testo in Processing

\begin{Shaded}
\begin{Highlighting}[]
\FunctionTok{textSize}\OperatorTok{(}\DecValTok{32}\OperatorTok{);} \CommentTok{//per la dimensione del testo}
\FunctionTok{text}\OperatorTok{(}\StringTok{"word"}\OperatorTok{,} \DecValTok{10}\OperatorTok{,} \DecValTok{30}\OperatorTok{);} \CommentTok{//così diciamo cosa vogliamo scrivere e indichiamo l\textquotesingle{}angolo in alto a sinitra}
\FunctionTok{fill}\OperatorTok{(}\DecValTok{0}\OperatorTok{,}\DecValTok{100}\OperatorTok{,}\DecValTok{150}\OperatorTok{);} \CommentTok{//con il fill gli diamo il colore}
\FunctionTok{text}\OperatorTok{(}\StringTok{"word"}\OperatorTok{,} \DecValTok{10}\OperatorTok{,} \DecValTok{60}\OperatorTok{);}
\FunctionTok{fill}\OperatorTok{(}\DecValTok{0}\OperatorTok{,}\DecValTok{100}\OperatorTok{,}\DecValTok{150}\OperatorTok{,}\DecValTok{50}\OperatorTok{);} \CommentTok{//ovviamente l\textquotesingle{}ultimo parametro è la trasparenza}
\FunctionTok{text}\OperatorTok{(}\StringTok{"word"}\OperatorTok{,} \DecValTok{10}\OperatorTok{,} \DecValTok{90}\OperatorTok{);}
\CommentTok{//possiamo anche cambiare il font}
\CommentTok{//createFont() }
\CommentTok{//loadFont() possiamo caricare un font creato con processing}

\NormalTok{PFont f}\OperatorTok{;} \CommentTok{//PFont sta per ProcessingFont, ed è una variabile dove mettiamo i font}
\NormalTok{f}\OperatorTok{=}\FunctionTok{createFont}\OperatorTok{(}\StringTok{"Arial"}\OperatorTok{,} \DecValTok{40}\OperatorTok{);} \CommentTok{//Come stringa gli passiamo una delle famiglie di font installate sul sistema, e passiamo anche la dimensione}
    \CommentTok{//E\textquotesingle{} meglio usare il tool "Create Font" di processing}
\FunctionTok{text}\OperatorTok{(}\StringTok{"word"}\OperatorTok{,}\DecValTok{10}\OperatorTok{,}\DecValTok{30}\OperatorTok{);} \CommentTok{//così diciamo cosa vogliamo scrivere e indichiamo l\textquotesingle{}angolo in alto a sinitra}
\FunctionTok{fill}\OperatorTok{(}\DecValTok{0}\OperatorTok{,}\DecValTok{100}\OperatorTok{,}\DecValTok{150}\OperatorTok{);} \CommentTok{//con il fill gli diamo il colore}
\FunctionTok{text}\OperatorTok{(}\StringTok{"word"}\OperatorTok{,}\DecValTok{10}\OperatorTok{,}\DecValTok{60}\OperatorTok{);}
\FunctionTok{fill}\OperatorTok{(}\DecValTok{0}\OperatorTok{,}\DecValTok{100}\OperatorTok{,}\DecValTok{150}\OperatorTok{,}\DecValTok{50}\OperatorTok{);} \CommentTok{//ovviamente l\textquotesingle{}ultimo parametro è la trasparenza}
\FunctionTok{text}\OperatorTok{(}\StringTok{"word"}\OperatorTok{,}\DecValTok{10}\OperatorTok{,}\DecValTok{90}\OperatorTok{);}
\NormalTok{PFont f}\OperatorTok{;} 
\NormalTok{f}\OperatorTok{=}\FunctionTok{loadFont}\OperatorTok{(}\StringTok{"Castellar{-}32.vlw"}\OperatorTok{);} 
\FunctionTok{textFont}\OperatorTok{(}\NormalTok{f}\OperatorTok{,}\DecValTok{28}\OperatorTok{);}

\FunctionTok{text}\OperatorTok{(}\StringTok{"word"}\OperatorTok{,}\DecValTok{10}\OperatorTok{,}\DecValTok{30}\OperatorTok{);} 
\FunctionTok{fill}\OperatorTok{(}\DecValTok{0}\OperatorTok{,}\DecValTok{100}\OperatorTok{,}\DecValTok{150}\OperatorTok{);} 
\FunctionTok{text}\OperatorTok{(}\StringTok{"word"}\OperatorTok{,}\DecValTok{10}\OperatorTok{,}\DecValTok{60}\OperatorTok{);}
\FunctionTok{fill}\OperatorTok{(}\DecValTok{0}\OperatorTok{,}\DecValTok{100}\OperatorTok{,}\DecValTok{150}\OperatorTok{,}\DecValTok{50}\OperatorTok{);} 
\FunctionTok{text}\OperatorTok{(}\StringTok{"word"}\OperatorTok{,}\DecValTok{10}\OperatorTok{,}\DecValTok{90}\OperatorTok{);}
\CommentTok{//Per creare un font, con l\textquotesingle{}utilizzo della funzione createFont(), basta andare in tools e selezionare il FontCreator, che crea il file .vlw  nella cartella data del nostro sketch}
\end{Highlighting}
\end{Shaded}

La prassi è:

\begin{itemize}
\item
  Andare su \textquotesingle Tool\textquotesingle{}
\item
  Selezionare \textquotesingle Create Font\textquotesingle{}
\item
  Selezionare un font
\item
  Copiare il nome del font + estensione (.vlw)
\item
  Nel codice creare una variabile di tipo
  \textquotesingle PFont\textquotesingle{}
\item
  Poi usare la funzione \textquotesingle loadFont()\textquotesingle{} e
  passare come argomento sottoforma di stringa il nome completo del font
  con l\textquotesingle estensione
\item
  Infine usare la funzione
  \textquotesingle textFont()\textquotesingle\textquotesingle{} dove
  inseriamo il nome della variabile PFont e la size
\end{itemize}

\begin{Shaded}
\begin{Highlighting}[]
\DataTypeTok{void} \FunctionTok{setup}\OperatorTok{()\{}
  \FunctionTok{size}\OperatorTok{(}\DecValTok{120}\OperatorTok{,}\DecValTok{120}\OperatorTok{);}
  \FunctionTok{textSize}\OperatorTok{(}\DecValTok{64}\OperatorTok{);}
  \FunctionTok{textAlign}\OperatorTok{(}\NormalTok{CENTER}\OperatorTok{);} \CommentTok{//Chiaramente serve per allineare il testo a sinitra, centro o a destra}
\OperatorTok{\}}

\DataTypeTok{void} \FunctionTok{draw}\OperatorTok{()\{}
  \FunctionTok{background}\OperatorTok{(}\DecValTok{0}\OperatorTok{);}
  \FunctionTok{text}\OperatorTok{(}\NormalTok{key}\OperatorTok{,} \DecValTok{60}\OperatorTok{,} \DecValTok{80}\OperatorTok{);} \CommentTok{//con key prende il carattere che pigiamo sulla tastiera}
\OperatorTok{\}}
\DataTypeTok{void} \FunctionTok{setup}\OperatorTok{()\{}
  \FunctionTok{size}\OperatorTok{(}\DecValTok{500}\OperatorTok{,}\DecValTok{500}\OperatorTok{);} \CommentTok{//Misura della canvas                                            }
  \FunctionTok{textSize}\OperatorTok{(}\DecValTok{42}\OperatorTok{);}  \CommentTok{//Dimensione del testo}
  \FunctionTok{fill}\OperatorTok{(}\DecValTok{0}\OperatorTok{);}       \CommentTok{//Colore del testo impostato uguale a nero}
  \BuiltInTok{String}\NormalTok{ x }\OperatorTok{=} \StringTok{"questa è una frase più lunga di una parola"}\OperatorTok{;}  \CommentTok{//Stringa denominata x che contiene la frase tra virgolette}
  \CommentTok{//text(x,10,30);  //Posizionamento alle coordinate 10,10 del punto in alto sinistra del rettangolo che contiene la stringa x}

  \CommentTok{//Per fare in modo che il testo non venga troncato ed esca dallo schermo:}
  \CommentTok{//Esiste una versione della funzione text che accetta un 4 e 5 parametro che ci permettono di definire lunghezza e altezza del rettangolo che contiene il testo}
  \FunctionTok{text}\OperatorTok{(}\NormalTok{x}\OperatorTok{,}\DecValTok{10}\OperatorTok{,}\DecValTok{30}\OperatorTok{,} \DecValTok{200}\OperatorTok{,}\DecValTok{300}\OperatorTok{);} \CommentTok{//In questo caso il rettangolo che contiene il testo ha dimensione 200X300, e il suo angolo in alto a sinistra è posizionato a 10,30}
\OperatorTok{\}}

\DataTypeTok{void} \FunctionTok{draw}\OperatorTok{()\{}
\OperatorTok{\}}
\end{Highlighting}
\end{Shaded}

\begin{Shaded}
\begin{Highlighting}[]
\CommentTok{/*Questo codice: trasforma la tua finestra di Processing in un mini{-}programma di videoscrittura, tipo il Blocco Note. Ti permette di digitare con la tastiera, cancellare gli errori e "salvare" la frase premendo Invio.*/}

\NormalTok{PFont f}\OperatorTok{;} \CommentTok{//Creiamo una variabile Processing Font \textquotesingle{}f\textquotesingle{} che contiene il font}
\BuiltInTok{String}\NormalTok{ scrivendo }\OperatorTok{=} \StringTok{""}\OperatorTok{;} \CommentTok{//Una strinfa vuota che serve a memorizzare, lettera per lettera, quello che stiamo digitando sil momento}
\BuiltInTok{String}\NormalTok{ salvata }\OperatorTok{=} \StringTok{""}\OperatorTok{;} \CommentTok{//Un\textquotesingle{}altra stringa vuota. Qui viene inserito il testo dopo aver premuto \textquotesingle{}invio\textquotesingle{}, per fingere di averlo salvato}

\DataTypeTok{void} \FunctionTok{setup}\OperatorTok{()\{}
  \FunctionTok{size}\OperatorTok{(}\DecValTok{500}\OperatorTok{,}\DecValTok{500}\OperatorTok{);} \CommentTok{//Misura della canvas                                            }
\NormalTok{  f}\OperatorTok{=}\FunctionTok{createFont}\OperatorTok{(}\StringTok{"Calibri"}\OperatorTok{,} \DecValTok{16}\OperatorTok{);} \CommentTok{//Carica il font "Calibri" con grandezza 16 e lo inserisce nella variabile f}
\OperatorTok{\}}

\DataTypeTok{void} \FunctionTok{draw}\OperatorTok{()\{}
  \FunctionTok{background}\OperatorTok{(}\DecValTok{255}\OperatorTok{);} \CommentTok{//Carica il font "Calibri" con grandezza 16 e lo infila nella scatola f}
  \DataTypeTok{int}\NormalTok{ ind }\OperatorTok{=} \DecValTok{25}\OperatorTok{;} \CommentTok{//Una variabile usata solo per comodità: è il margine sinistro (x)}
  \FunctionTok{textFont}\OperatorTok{(}\NormalTok{f}\OperatorTok{);} \CommentTok{//Diciamo a Processing: "Usa il font che ho caricato prima"}
  \FunctionTok{fill}\OperatorTok{(}\DecValTok{0}\OperatorTok{);} \CommentTok{//Colore nero per il testo}
    
  \CommentTok{//Scriviamo sullo schermo usando: \textquotesingle{}text("Frase", posizioneX, posizioneY)\textquotesingle{}}
  \FunctionTok{text}\OperatorTok{(}\StringTok{"Click sulla finestra.}\SpecialCharTok{\textbackslash{}n}\StringTok{Premi invio per salvare. "}\OperatorTok{,}\NormalTok{ ind}\OperatorTok{,} \DecValTok{40}\OperatorTok{);}
  \FunctionTok{text}\OperatorTok{(}\StringTok{"Input: "}\OperatorTok{+}\NormalTok{scrivendo}\OperatorTok{,}\NormalTok{ ind}\OperatorTok{,} \DecValTok{190}\OperatorTok{);}
  \FunctionTok{text}\OperatorTok{(}\StringTok{"Testo salvato: "}\OperatorTok{+}\NormalTok{salvata}\OperatorTok{,}\NormalTok{ ind}\OperatorTok{,} \DecValTok{230}\OperatorTok{);}
\OperatorTok{\}}

\CommentTok{//Questa funzione scatta in automatico solo e unicamente quando premi un tasto sulla tastiera. Processing salva il tasto che hai premuto in una variabile di sistema chiamata \textquotesingle{}key\textquotesingle{}}
\DataTypeTok{void} \FunctionTok{keyPressed}\OperatorTok{()\{}
    \CommentTok{// CASO 1: L\textquotesingle{}utente preme INVIO}
    \ControlFlowTok{if}\OperatorTok{(}\NormalTok{key}\OperatorTok{==}\CharTok{\textquotesingle{}\textbackslash{}n\textquotesingle{}}\OperatorTok{)\{}
\NormalTok{      salvata}\OperatorTok{=}\NormalTok{scrivendo}\OperatorTok{;} \CommentTok{//Copiamo la frase attuale dentro la variabile "salvata"}
\NormalTok{      scrivendo}\OperatorTok{=}\StringTok{""}\OperatorTok{;} 	 \CommentTok{//Svuotiamo la riga di digitazione per ricominciare da capo!}
    \OperatorTok{\}}
    \CommentTok{// CASO 2: L\textquotesingle{}utente preme CANCELLA (Backspace)}
    \ControlFlowTok{else} \ControlFlowTok{if}\OperatorTok{(}\NormalTok{key}\OperatorTok{==}\NormalTok{BACKSPACE }\OperatorTok{\&\&}\NormalTok{ scrivendo}\OperatorTok{.}\FunctionTok{length}\OperatorTok{()} \OperatorTok{\textgreater{}} \DecValTok{0}\OperatorTok{)\{}
\CommentTok{//scrivendo.length() \textgreater{} 0: Questo controllo serve a evitare che il programma crashi. Gli diciamo: "Cancella l\textquotesingle{}ultima lettera solo se c\textquotesingle{}è almeno una lettera scritta! Se la frase è già vuota, non fare niente".}
        
      \CommentTok{// Substring "taglia" la parola: prende tutto dalla prima lettera (0) fino alla penultima.}
\CommentTok{//substring(inizio, fine): È un paio di forbici digitali. Se la stringa è "CIAO" (lunghezza 4), substring(0, 3) taglierà dalla lettera alla posizione 0 (la C) fino alla 3 esclusa. Risultato? "CIA", ha tolto la O!}
\NormalTok{      scrivendo}\OperatorTok{=}\NormalTok{scrivendo}\OperatorTok{.}\FunctionTok{substring}\OperatorTok{(}\DecValTok{0}\OperatorTok{,}\NormalTok{scrivendo}\OperatorTok{.}\FunctionTok{length}\OperatorTok{(){-}}\DecValTok{1}\OperatorTok{);}
    \OperatorTok{\}}
    \CommentTok{// CASO 3: L\textquotesingle{}utente preme una qualsiasi altra lettera}
    \ControlFlowTok{else}\OperatorTok{\{}
\NormalTok{      scrivendo}\OperatorTok{+=}\NormalTok{key}\OperatorTok{;} \CommentTok{//Aggiunge la lettera digitata alla fine della frase, quindi una dopo l\textquotesingle{}altra}
    \OperatorTok{\}}
\OperatorTok{\}}
\end{Highlighting}
\end{Shaded}

Chiariamo la sintassi che contiene i cosiddetti dot
\textquotesingle.\textquotesingle{}\\
Tipica di Java.

In Java, una String (una parola) non è solo un contenitore di lettere: è
un Oggetto. Un oggetto contiene l\textquotesingle informazione (la
parola stessa), possiede anche gli altri strumenti (chiamati metodi) per
farsi manipolare.

Il punto (\texttt{.}) è il modo in cui tu "premi il pulsante" per
attivare uno di quegli strumenti su quell\textquotesingle oggetto
specifico.

\begin{itemize}
\item
  \texttt{scrivendo.length()} significa: \emph{"Ehi stringa
  \texttt{scrivendo}, dimmi la tua lunghezza!"}.
\item
  \texttt{scrivendo.substring(...)} significa: \emph{"Ehi stringa
  \texttt{scrivendo}, usa il tuo strumento interno per ritagliarti!"}.
\end{itemize}

Come funziona substring() in generale?\\
La parola \texttt{substring} significa letteralmente "sotto-stringa"
(una porzione della stringa). Serve a estrarre un pezzetto di testo da
una parola più lunga.

Per usarlo devi ricordarti la solita regola: \textbf{i computer iniziano
a contare da 0, non da 1.} Immagina la stringa \texttt{"CIAO"}:

\begin{itemize}
\item
  La lettera \textbf{C} è nella posizione \textbf{0}
\item
  La lettera \textbf{I} è nella posizione \textbf{1}
\item
  La lettera \textbf{A} è nella posizione \textbf{2}
\item
  La lettera \textbf{O} è nella posizione \textbf{3}
\end{itemize}

La sintassi di base è questa:
\texttt{parola.substring(indiceDiPartenza,\ indiceDiArrivo);}

\textbf{⚠️ ATTENZIONE ALLA REGOLA TRABOCCHETTO:}

\begin{itemize}
\item
  L\textquotesingle{}\texttt{indiceDiPartenza} è \textbf{incluso} nel
  taglio.
\item
  L\textquotesingle{}\texttt{indiceDiArrivo} è \textbf{escluso} dal
  taglio. Il computer si ferma \emph{un attimo prima} di quel vagone.
\end{itemize}

\textbf{Esempi pratici con "CIAO":}

\begin{itemize}
\item
  \texttt{"CIAO".substring(0,\ 2)} ➔ Taglia da 0 a 1 (il 2 è escluso).
  Risultato: \textbf{"CI"}.
\item
  \texttt{"CIAO".substring(1,\ 4)} ➔ Taglia da 1 a 3 (il 4 è escluso).
  Risultato: \textbf{"IAO"}
\end{itemize}

\begin{center}\rule{0.5\linewidth}{0.5pt}\end{center}

Adesso abbiamo creato un motore fisico in miniatura. Stiamo simulando la
gravità e l\textquotesingle attrito (la perdita di energia) nel mondo
reale.

\begin{Shaded}
\begin{Highlighting}[]
\DataTypeTok{float}\NormalTok{ x }\OperatorTok{=} \DecValTok{240}\OperatorTok{;} \CommentTok{//Coordinata x dell\textquotesingle{}angolo in alto a sinistra del rettangolo}
\DataTypeTok{float}\NormalTok{ y }\OperatorTok{=} \DecValTok{0}\OperatorTok{;} \CommentTok{//Coordinata y dell\textquotesingle{}angolo in alto a sinistra del rettangolo}
\DataTypeTok{float}\NormalTok{ speed }\OperatorTok{=} \DecValTok{0}\OperatorTok{;} \CommentTok{//La velocità iniziale è 0. Immagina di tenere l\textquotesingle{}oggetto in mano a mezz\textquotesingle{}aria prima di lasciarlo andare. Non sta ancora cadendo.}
\DataTypeTok{float}\NormalTok{ gravity }\OperatorTok{=} \FloatTok{0.1}\OperatorTok{;} \CommentTok{//La "forza" che spinge verso il basso. Vale 0.1 e agirà costantemente.}

\DataTypeTok{void} \FunctionTok{setup}\OperatorTok{()\{}
    \FunctionTok{size}\OperatorTok{(}\DecValTok{500}\OperatorTok{,}\DecValTok{250}\OperatorTok{);}
\OperatorTok{\}}

\DataTypeTok{void} \FunctionTok{draw}\OperatorTok{()\{}  
\CommentTok{//Qui niente di nuovo, pulisci lo schermo e disegni il rettangolino grigio. Nota l\textquotesingle{}ottima scelta di usare rectMode(CENTER): in questo modo la variabile y indica esattamente il centro del rettangolo, e non il suo angolo in alto a sinistra.}
    \FunctionTok{background}\OperatorTok{(}\DecValTok{255}\OperatorTok{);}
    \FunctionTok{fill}\OperatorTok{(}\DecValTok{175}\OperatorTok{);}
    \FunctionTok{stroke}\OperatorTok{(}\DecValTok{0}\OperatorTok{);}
    \FunctionTok{rectMode}\OperatorTok{(}\NormalTok{CENTER}\OperatorTok{);} 
    \FunctionTok{rect}\OperatorTok{(}\NormalTok{x}\OperatorTok{,}\NormalTok{y}\OperatorTok{,}\DecValTok{10}\OperatorTok{,}\DecValTok{10}\OperatorTok{);}
    
    \CommentTok{//La gravità in azione}
\NormalTok{    y}\OperatorTok{+=}\NormalTok{speed}\OperatorTok{;} \CommentTok{//La posizione Y cambia in base alla velocità. (Se la velocità è 2, il quadrato scende di 2 pixel).}
\NormalTok{    speed}\OperatorTok{+=}\NormalTok{gravity}\OperatorTok{;} \CommentTok{//Al primo giro la velocità è 0. Poi diventa 0.1, poi 0.2, poi 0.3... Più scende, più la sua velocità aumenta.}
    
    \CommentTok{//L\textquotesingle{}urto ed il rimbalzo}
    \ControlFlowTok{if}\OperatorTok{(}\NormalTok{y }\OperatorTok{\textgreater{}}\NormalTok{ height}\OperatorTok{)\{} \CommentTok{//In questo caso sta urtando con la parte bassa della canvas. Quindi deve rimbalzare, e per rimbalzare dobbiamo cambiare il senso della sua velocità.}
\NormalTok{      speed }\OperatorTok{=}\NormalTok{ speed }\OperatorTok{*} \OperatorTok{{-}}\FloatTok{0.95}\OperatorTok{;} \CommentTok{//Il segno Meno ({-}) inverte la direzione. Se la velocità era, per esempio, +10 (verso il basso), moltiplicandola per un numero negativo diventa {-}9.5 (verso l\textquotesingle{}alto). Il 0.95 simula l\textquotesingle{}attrito (o la dispersione di energia). Significa: "mantieni solo il 95\% della velocità che avevi al momento dello schianto". Quel 5\% perso fa sì che ogni rimbalzo sia un po\textquotesingle{} più basso del precedente, finché non si ferma.}
\NormalTok{      y }\OperatorTok{=}\NormalTok{ height}\OperatorTok{;} \CommentTok{//Quando il computer fa i calcoli, a volte il quadrato potrebbe trovarsi a Y = 252, "sfondando" un po\textquotesingle{} il pavimento (che è a 250). Questa riga lo "riacchiappa" e lo riposiziona esattamente a filo col pavimento, evitando che si incastri o faccia scatti strani.}
    \OperatorTok{\}}
\OperatorTok{\}}
\end{Highlighting}
\end{Shaded}

\begin{Shaded}
\begin{Highlighting}[]
\DataTypeTok{int}\NormalTok{ a }\OperatorTok{=} \DecValTok{0}\OperatorTok{;}

\DataTypeTok{void} \FunctionTok{setup}\OperatorTok{()\{}
  \FunctionTok{size}\OperatorTok{(}\DecValTok{500}\OperatorTok{,}\DecValTok{500}\OperatorTok{);}
    \CommentTok{//rectMode(CENTER); //Se vogliamo fare in modo che il rettangolo giri su se stesso}
\OperatorTok{\}}

\DataTypeTok{void} \FunctionTok{draw}\OperatorTok{()\{}
  \FunctionTok{background}\OperatorTok{(}\DecValTok{0}\OperatorTok{);}
  \FunctionTok{translate}\OperatorTok{(}\NormalTok{width}\OperatorTok{/}\DecValTok{2}\OperatorTok{,}\NormalTok{height}\OperatorTok{/}\DecValTok{2}\OperatorTok{);}
  \CommentTok{//Il foglio è sempre largo 500 e alto 500, ma ora le coordinate per muoverti sono cambiate: Il centro esatto è diventato (0, 0). L\textquotesingle{}angolo in alto a sinistra ora è a coordinate ({-}250, {-}250). L\textquotesingle{}angolo in basso a destra ora è a coordinate (+250, +250).}
  \CommentTok{//Se tu mettessi translate nel setup(), Processing sposterebbe l\textquotesingle{}origine al centro una volta sola all\textquotesingle{}avvio. Ma appena partirebbe il primo draw(), il sistema se ne dimenticherebbe all\textquotesingle{}istante, resettando la griglia in alto a sinistra. Ecco perché devi "ricordargli" di spostare il righello al centro 60 volte al secondo!}
  \FunctionTok{rotate}\OperatorTok{(}\FunctionTok{radians}\OperatorTok{(}\NormalTok{a}\OperatorTok{));} \CommentTok{//Questa è la funzione per la rotazione.}
  \FunctionTok{rect}\OperatorTok{(}\DecValTok{0}\OperatorTok{,}\DecValTok{0}\OperatorTok{,} \DecValTok{50}\OperatorTok{,}\DecValTok{50}\OperatorTok{);}
\NormalTok{  a}\OperatorTok{++;}
\OperatorTok{\}}
\end{Highlighting}
\end{Shaded}

\begin{Shaded}
\begin{Highlighting}[]
\DataTypeTok{void} \FunctionTok{setup}\OperatorTok{()\{}
    \FunctionTok{size}\OperatorTok{(}\DecValTok{500}\OperatorTok{,}\DecValTok{500}\OperatorTok{);}
\OperatorTok{\}}

\DataTypeTok{void} \FunctionTok{draw}\OperatorTok{()\{}
    \FunctionTok{background}\OperatorTok{(}\DecValTok{0}\OperatorTok{);}
    \FunctionTok{translate}\OperatorTok{(}\NormalTok{width}\OperatorTok{/}\DecValTok{2}\OperatorTok{,}\NormalTok{ height}\OperatorTok{/}\DecValTok{2}\OperatorTok{);} \CommentTok{//Trasliamo il nostro sistema di riferimento. Adesso width/2 ed height/2 sono i nostri nuovi 0,0}
    \CommentTok{/*}
\CommentTok{    millis(): È un cronometro interno di Processing. Ti dice quanti millisecondi sono passati da quando hai premuto Play.}
\CommentTok{    * 0.001: Trasforma i millisecondi in secondi (es. 1500 millisecondi diventano 1.5 secondi).}
\CommentTok{    TWO\_PI: È il valore di 360 gradi scritti in radianti (un giro completo).}

\CommentTok{    Significa "Fai un giro completo esatto ogni secondo"}
\CommentTok{    Dividendo il tutto per 10 (come nel primo rotate), significa: "Fai un giro completo ogni 10 secondi" (molto più lento!).}
\CommentTok{    */}
    
    \CommentTok{//Con millis * 0.001 stiamo convertendo i millisecodi in secondi. Con TWO\_PI/10 stiamo dicendo "Ogni secondo fa un decimo di giro", quindi un giro completo ogni 10 secondi}
    \FunctionTok{rotate}\OperatorTok{(}\FunctionTok{millis}\OperatorTok{()*}\FloatTok{0.001}\OperatorTok{*}\NormalTok{TWO\_PI}\OperatorTok{/}\DecValTok{10}\OperatorTok{);} \CommentTok{//Rivoluzione attorno al Sole}
    \FunctionTok{translate}\OperatorTok{(}\DecValTok{80}\OperatorTok{,}\DecValTok{0}\OperatorTok{);} \CommentTok{//Spostati di 80px a destra, Quando usi rotate() e subito dopo translate(X, 0), non ti sposti a destra sullo schermo, ma cammini in avanti lungo il tuo asse appena ruotato. In questo caso, X diventa letteralmente il RAGGIO della tua orbita.}
    \FunctionTok{rotate}\OperatorTok{(}\FunctionTok{millis}\OperatorTok{()*}\FloatTok{0.001}\OperatorTok{*}\NormalTok{TWO\_PI}\OperatorTok{);} \CommentTok{//Questa fa girare il pianeta su se stesso mentre sta compiendo l\textquotesingle{}orbita}

    \CommentTok{//Gestione della ellisse (Spessore, colore e posizione [posizione ovviamente con il nuovo sistema di riferimento])}
    \FunctionTok{fill}\OperatorTok{(}\DecValTok{255}\OperatorTok{);}
    \FunctionTok{noStroke}\OperatorTok{();}
    \FunctionTok{ellipseMode}\OperatorTok{(}\NormalTok{CENTER}\OperatorTok{);}
    \FunctionTok{ellipse}\OperatorTok{(}\DecValTok{0}\OperatorTok{,}\DecValTok{0}\OperatorTok{,}\DecValTok{40}\OperatorTok{,}\DecValTok{40}\OperatorTok{);}
    
    \CommentTok{//Gestione della linea (Spessore, colore e posizione [posizione ovviamente con il nuovo sistema di riferimento])}
    \FunctionTok{strokeWeight}\OperatorTok{(}\DecValTok{5}\OperatorTok{);}
    \FunctionTok{stroke}\OperatorTok{(}\DecValTok{255}\OperatorTok{,}\DecValTok{0}\OperatorTok{,}\DecValTok{0}\OperatorTok{);}
    \FunctionTok{line}\OperatorTok{(}\DecValTok{0}\OperatorTok{,}\DecValTok{0}\OperatorTok{,}\DecValTok{40}\OperatorTok{,}\DecValTok{0}\OperatorTok{);}
\OperatorTok{\}}
\end{Highlighting}
\end{Shaded}

\begin{Shaded}
\begin{Highlighting}[]
\DataTypeTok{int}\NormalTok{ x }\OperatorTok{=} \DecValTok{0}\OperatorTok{;}
\DataTypeTok{void} \FunctionTok{setup}\OperatorTok{()\{}
    \FunctionTok{size}\OperatorTok{(}\DecValTok{500}\OperatorTok{,}\DecValTok{500}\OperatorTok{);}
    \FunctionTok{rectMode}\OperatorTok{(}\NormalTok{CENTER}\OperatorTok{);}
\OperatorTok{\}}

\DataTypeTok{void} \FunctionTok{draw}\OperatorTok{()\{}
\CommentTok{//Essendo che rotate e translate sono due operazioni "Matriciali", possiamo creare dei movimenti indipendenti, e quindi evitare che più oggetti nello schermo seguano lo stesso movimento}
\CommentTok{//Questo viene fatto usano pushMatrix() e popMatrix(). Li usiamo come due fette di pane, che racchiudono nel panino tutto quello che riguarda oggetti e movimento}

\FunctionTok{background}\OperatorTok{(}\DecValTok{255}\OperatorTok{);}
\CommentTok{//Facciamo un cubo verde con un movimento proprio ed una velocità propria}
    \FunctionTok{pushMatrix}\OperatorTok{();}
    \FunctionTok{translate}\OperatorTok{(}\NormalTok{x}\OperatorTok{+}\DecValTok{300}\OperatorTok{,}\DecValTok{200}\OperatorTok{);} \CommentTok{//Occhio alla x, Se la togliessimo il cubo non si muoverebbe a destra}
    \FunctionTok{rotate}\OperatorTok{(}\FunctionTok{radians}\OperatorTok{(}\NormalTok{frameCount}\OperatorTok{));}
    \FunctionTok{fill}\OperatorTok{(}\DecValTok{0}\OperatorTok{,}\DecValTok{255}\OperatorTok{,}\DecValTok{0}\OperatorTok{);}
    \FunctionTok{rect}\OperatorTok{(}\DecValTok{0}\OperatorTok{,}\DecValTok{0}\OperatorTok{,}\DecValTok{40}\OperatorTok{,}\DecValTok{40}\OperatorTok{);}
    \FunctionTok{popMatrix}\OperatorTok{();}

\CommentTok{//Facciamo un cubo rosso con un movimento proprio ed una velocità propria}
    \FunctionTok{pushMatrix}\OperatorTok{();}
    \FunctionTok{translate}\OperatorTok{(}\NormalTok{x}\OperatorTok{+}\DecValTok{100}\OperatorTok{,}\DecValTok{200}\OperatorTok{);}  \CommentTok{//Occhio alla x, Se la togliessimo il cubo non si muoverebbe a destra}
    \FunctionTok{rotate}\OperatorTok{(}\FunctionTok{radians}\OperatorTok{(}\NormalTok{frameCount}\OperatorTok{)*}\DecValTok{2}\OperatorTok{);}
    \FunctionTok{fill}\OperatorTok{(}\DecValTok{255}\OperatorTok{,}\DecValTok{0}\OperatorTok{,}\DecValTok{0}\OperatorTok{);}
    \FunctionTok{rect}\OperatorTok{(}\DecValTok{0}\OperatorTok{,}\DecValTok{0}\OperatorTok{,}\DecValTok{40}\OperatorTok{,}\DecValTok{40}\OperatorTok{);}
    \FunctionTok{popMatrix}\OperatorTok{();}

 \CommentTok{//La x, inoltre incrementata sempre di 1, serve a far spostare verso destra i cubi}
\NormalTok{    x}\OperatorTok{++;}
\OperatorTok{\}}
\end{Highlighting}
\end{Shaded}


\end{document}
